%% kapa: exam=1 round=19 year=2008 period=1 form=A subject=appraisal_law subject_ko=부동산관계법규 seq=29 number=69
%% source: Q-Net 감정평가사 기출문제 (cst00309) artlSeq=1001962 file=extracted/1008/1001962/2008년 제19회 감정평가사 1차시험 기출문제/1교시.hwp
%% conversion: kordoc 4.24.51 -> tools/gen1.py (auto); review=proofread-sonnet; flags=none
\begin{question}{69}
건축법령상 1 이상의 필지의 일부를 하나의 대지로 할 수 있는 경우가 아닌 것은 ?
\begin{choices}
  \item 1 이상의 필지의 일부에 대하여 도시계획시설이 결정․고시가 있는 부분의 토지
  \item 1 이상의 필지의 일부에 대하여 농지법에 의한 농지전용허가를 받은 부분의 토지
  \item 1 이상의 필지의 일부에 대하여 산지관리법에 의한 산지전용허가를 받은 부분의 토지
  \item 1 이상의 필지의 일부에 대하여 국토의 계획 및 이용에 관한 법률에 의한 개발행위허가를 받은 부분의 토지
  \item 1 이상의 필지의 일부에 대하여 하천법에 의한 하천점용허가를 받은 부분의 토지
\end{choices}
\end{question}
\answer{5}
